\section{Optimality of degree five}\label{sec:minimality}

The construction began with a quintic because the Berend--Bilu example is
a product of degrees three and two.  It is natural to ask whether the same
fixed-map phenomenon could have been found in smaller degree.  Here the
arithmetic answers before the geometry does: a nontrivially intersective
polynomial cannot have degree below five.  The degree in which the example
first appears is, pleasantly, also the first degree in which it could appear.

Define \(d_{\HP}\) to be the least geometric degree of a polynomial Keller
map over \(\Q\) having a finite regular fiber that satisfies
\eqref{eq:hasse-condition}.

\begin{theorem}\label{thm:minimality}
The minimal Hasse-failing geometric degree is
\[
 \boxed{d_{\HP}=5.}
 \tag{5.1}
\]
\end{theorem}

\begin{proof}
Let \(F:\A^n_{\Q}\to\A^n_{\Q}\) be a Keller map of geometric degree \(d\),
and let \(X=F^{-1}(y)\) be a finite regular fiber.  Because \(F\) is
\'etale, each geometric point of \(X\) extends, after shrinking around
\(y\), to a distinct finite \'etale branch.  The geometric generic fiber
has \(d\) points, so
\[
 \operatorname{rank}_{\Q}X\le d.
 \tag{5.2}
\]

Suppose \(X\) is regular, everywhere locally soluble, and has no rational
point.  Choose an integral primitive element for its finite \'etale
coordinate algebra.  Thus
\[
 X\simeq\Spec\Q[T]/(P)
 \tag{5.3}
\]
for a monic squarefree polynomial \(P\in\Z[T]\) of degree
\(\operatorname{rank}_{\Q}X\).  A \(\Q_p\)-root of the monic polynomial is
integral.  Roots over every \(\Q_p\) therefore give roots modulo every prime
power, and the Chinese remainder theorem gives a root modulo every positive
integer.  The absence of a \(\Q\)-point says that \(P\) has no rational
root.

For completeness, the low-degree consequence of Berend and Bilu's
criterion can be seen directly from the possible irreducible-factor
degrees \cite{BerendBilu}.  With no rational root, a polynomial of total
degree at most four is either irreducible of degree \(2\), \(3\), or \(4\),
or is a product of two irreducible quadratics.  In the irreducible case a
derangement in the transitive Galois action gives a Frobenius class with no
fixed root.  In the two-quadratic case, if the quadratic fields agree then
an inert Frobenius works for both factors; if they are distinct, their
biquadratic compositum contains an element nontrivial in both quadratic
quotients.  The Berend--Bilu fixed-point criterion (equivalently,
Chebotarev applied to these classes) then supplies a prime modulo which
there is no root.  Hence no polynomial of degree at most four can have a
root modulo every positive integer without already having a rational root.
Thus
\[
 5\le\operatorname{rank}_{\Q}X\le d.
 \tag{5.4}
\]
On the other hand, Theorem~\ref{thm:main} gives a full Hasse-failing fiber
of a degree-five Keller map.  Hence \(d_{\HP}=5\).
\end{proof}

\begin{remark}
The lower bound is arithmetic and does not depend on the quadratic-gauge
construction.  The upper bound is stronger than mere realizability: one
fixed degree-five map attains the minimum at infinitely many targets.
\end{remark}
